% TOP_PROBLEM: {"id": 1308, "title": "Andrews–Curtis Conjecture", "category_id": 4, "category": {"id": 4, "name": "algebra", "display_name": "Algebra", "description": "Group theory, ring theory, field theory, and algebraic structures.", "slug": "algebra", "order_index": 4, "created_at": "2026-07-31T15:26:25.671Z"}, "status": "open"}
% FIRST_POSTED: null
% ENTRY_KIND: statement_only
% Imported from: batch05.tex; UnsolvedMath ID 1308
\documentclass[11pt]{article}
\usepackage[margin=25mm]{geometry}
\usepackage{fontspec}
\setmainfont{Latin Modern Roman}
\usepackage{amsmath,amssymb,mathtools}
\usepackage{unicode-math}
\setmathfont{Latin Modern Math}
\usepackage{xurl,hyperref}
\hypersetup{colorlinks=true,urlcolor=blue}
\setcounter{secnumdepth}{0}
\setlength{\parindent}{0pt}
\setlength{\parskip}{5pt}
\setlength{\emergencystretch}{3em}
\newcommand{\sref}[2]{\hyperlink{source:#1:#2}{[S#2]}}
\newcommand{\cataloguescope}[1]{\paragraph{Recorded target.}#1}
\begin{document}
% UnsolvedMath ID 1308; source catalogue code ALG-012
\setcounter{section}{1307}
\section{Andrews–Curtis Conjecture}\label{1308}
{\small\sffamily UnsolvedMath ID 1308 · Algebra}
\subsection{Definitions and mathematical statement}

\paragraph{Shared notation.}$\mathbb N=\{1,2,\ldots\}$, and $\mathbb Z,\mathbb Q,\mathbb R,\mathbb C$ denote the integers, rationals, reals and complex numbers. Any other conventions are specified locally.


Let $F_n=\langle x_1,\ldots,x_n\rangle$ be the free group on $n\geq1$ generators. A presentation $\langle x_1,\ldots,x_n\mid r_1,\ldots,r_n\rangle$ is balanced and presents the trivial group exactly when the normal closure $\langle\!\langle r_1,\ldots,r_n\rangle\!\rangle$ equals $F_n$. The allowed Andrews--Curtis moves on an ordered relator tuple are: replace $r_i$ by $r_i^{-1}$; replace $r_i$ by $r_ir_j$ for $j\ne i$; or replace $r_i$ by $wr_iw^{-1}$ for any $w\in F_n$. Moves and their inverses may be repeated finitely many times; interchanging two relators is also permitted, equivalently obtainable using these moves. Free reduction does not change the represented group word. Generators and the number of relators remain fixed.
\paragraph{Statement recorded in the supplied source.}
For every $n\geq1$ and every tuple $(r_1,\ldots,r_n)$ normally generating $F_n$, can these moves transform it into $(x_1,\ldots,x_n)$? Equivalently, is every balanced presentation of the trivial group Andrews--Curtis equivalent to the same-rank trivial presentation? \sref{1308}{2}

\subsection{Short English statement}
Can every balanced presentation of the trivial group be simplified to the standard trivial presentation using inversion, relator multiplication and conjugation, without adding generators?
\subsection{Sources}
\begin{description}
\item[\hypertarget{source:1308:1}{S1}] ULAM.AI and the UnsolvedMath Contributors, UnsolvedMath: A Curated Collection of Open Mathematics Problems, record 1308 (ALG-012). CC BY 4.0. \url{https://huggingface.co/datasets/ulamai/UnsolvedMath}.
\item[\hypertarget{source:1308:2}{S2}] Alexei Lisitsa, \emph{Stable Andrews--Curtis trivialization of AK(3) revisited. A case study using automated deduction}, arXiv:2501.18601v1 (2025), Section 1, moves AC1--AC3 and Conjecture 1. \url{https://arxiv.org/pdf/2501.18601}.
\end{description}
\label{end:1308}
\end{document}
